% Source: Solutions/Ansh-PS1-Q1-5.pdf, pp. 2-3 (typed). Restated problem statement omitted.
\submittedby{Ansh}

\textbf{Answer.} No, $\tau_\infty$ is not a topology on $X$. It violates both the first axiom
($X\in\tau_\infty$) and the second axiom (closure under arbitrary unions; in fact, it even fails
closure under finite unions).

\textbf{Justification.} For $\tau_\infty$ to be a topology on $X$, it must satisfy the three
standard topological axioms:
\begin{itemize}
  \item[(T1)] $\emptyset\in\tau_\infty$ and $X\in\tau_\infty$;
  \item[(T2)] $\tau_\infty$ is closed under arbitrary unions;
  \item[(T3)] $\tau_\infty$ is closed under finite intersections.
\end{itemize}
We examine each axiom in turn:

\textbf{1. Violation of Axiom (T1) ($X\notin\tau_\infty$).} By explicit definition,
$\emptyset\in\tau_\infty$. However, consider the whole space $U=X$. Since $X$ is an infinite
set, $X\neq\emptyset$. Furthermore, the complement of $X$ in $X$ is
\[ X\setminus X=\emptyset, \]
which has cardinality $0$ and is therefore a finite set, not an infinite set. Because $U=X$
neither equals $\emptyset$ nor has an infinite complement in $X$, we have
\[ X\notin\tau_\infty. \]
Thus, Axiom (T1) is violated.

\textbf{2. Violation of Axiom (T2) (Failure of closure under unions).} Even if one were to
adjoin $X$ to the collection, $\tau_\infty$ fails to be closed under arbitrary unions---and
even fails under the union of two sets:
\begin{itemize}
  \item \textbf{Counterexample via arbitrary unions of singletons:} For any element $x\in X$,
        consider the singleton $U_x=\{x\}$. Since $X$ is infinite, removing a single element
        leaves an infinite set, so $X\setminus\{x\}$ is infinite and $\{x\}\in\tau_\infty$
        for every $x\in X$. Fix a point $x_0\in X$ and consider the union
        \[ U:=\bigcup_{x\in X\setminus\{x_0\}}\{x\}=X\setminus\{x_0\}. \]
        Since $X$ is infinite, $U=X\setminus\{x_0\}$ is non-empty (and $U\neq X$). Yet its
        complement is $X\setminus U=\{x_0\}$, which is a finite set of cardinality $1$.
        Therefore, $U\notin\tau_\infty$.
  \item \textbf{Counterexample via a finite union of two sets:} Since $X$ is infinite, $X$
        contains a countably infinite subset $C=\{x_1,x_2,x_3,\dots\}$ with $x_i\neq x_j$ for
        $i\neq j$. Partition $C$ into the set of even-indexed elements
        $E=\{x_2,x_4,x_6,\dots\}$ and odd-indexed elements $O=\{x_1,x_3,x_5,\dots\}$, both of
        which are infinite. Define
        \[ U_1:=X\setminus(E\cup\{x_1\})\quad\text{and}\quad U_2:=X\setminus O. \]
        Then $X\setminus U_1=E\cup\{x_1\}$ is infinite, so $U_1\in\tau_\infty$; similarly,
        $X\setminus U_2=O$ is infinite, so $U_2\in\tau_\infty$. However, by De Morgan's laws,
        \[ U_1\cup U_2=X\setminus\big((E\cup\{x_1\})\cap O\big)=X\setminus\{x_1\}. \]
        Since $U_1\cup U_2=X\setminus\{x_1\}\neq\emptyset$ and its complement
        $X\setminus(U_1\cup U_2)=\{x_1\}$ is finite (of size $1$), we have
        $U_1\cup U_2\notin\tau_\infty$. (Moreover, taking $V_1=X\setminus E\in\tau_\infty$ and
        $V_2=X\setminus O\in\tau_\infty$ gives $V_1\cup V_2=X\notin\tau_\infty$.)
\end{itemize}

\textbf{3. Remark on Axiom (T3) (Closure under finite intersections holds).} For completeness,
note that if $U_1,U_2\in\tau_\infty$, then either one of them is $\emptyset$ (so
$U_1\cap U_2=\emptyset\in\tau_\infty$) or both $X\setminus U_1$ and $X\setminus U_2$ are
infinite, in which case $X\setminus(U_1\cap U_2)=(X\setminus U_1)\cup(X\setminus U_2)$ contains
$X\setminus U_1$ and is therefore infinite, so $U_1\cap U_2\in\tau_\infty$. Nevertheless,
because Axioms (T1) and (T2) fail, $\tau_\infty$ is not a topology on $X$.
